\begin{theorem}[No end-compatible whole-tile square with nontrivial kernel]\label{thm:whole-tile-no-go}
Let \eqref{eq:resolved-factorization} be a canonical resolved system with the annular-tower typing of Definition~\ref{def:canonical-resolved-system}, and let \eqref{eq:kl-tile-map} be an actual KL tile system as in Definition~\ref{def:kl-tile-system}.  If a whole-tile end-compatible comparison \((\Phi,B)\) in the sense of Definition~\ref{def:whole-tile-comparison} exists, then
\begin{equation}\label{eq:forced-trivial-kernel}
N=\{1\}.
\end{equation}
Equivalently, if \(N\ne\{1\}\), no pair \((\Phi,B)\) satisfies clauses (1)--(5) of Definition~\ref{def:whole-tile-comparison}.
\end{theorem}

\begin{proof}
Assume \(N\ne\{1\}\).  By Lemma~\ref{lem:wrapping},
\[
\zeta_+\in\Lambda(N).
\]
Choose \(z_0\in\mathcal E\).  By the definition of the limit set, there are distinct \(\gamma_n\in N\) such that
\begin{equation}\label{eq:deck-seq}
\gamma_nz_0\longrightarrow\zeta_+.
\end{equation}
For arbitrary \(z\in\mathcal E\), every \(\gamma_n\) is a hyperbolic isometry, so
\begin{equation}\label{eq:isometry}
d_{\mathcal E}(\gamma_nz,\gamma_nz_0)=d_{\mathcal E}(z,z_0).
\end{equation}
Bounded hyperbolic distance and \eqref{eq:deck-seq} imply
\begin{equation}\label{eq:all-z}
\gamma_nz\longrightarrow\zeta_+\qquad\text{for every }z\in\mathcal E,
\end{equation}
since hyperbolic disks of a fixed radius centered at points converging to a point of \(\partial\mathbb D\) have Euclidean diameters tending to \(0\).

Because \(N=\operatorname{Deck}(q)\),
\begin{equation}\label{eq:deck-invariance}
L(\gamma_nz)=e\,q(\gamma_nz)=e\,q(z)=L(z).
\end{equation}
Apply the composite square \eqref{eq:comparison-square}:
\begin{equation}\label{eq:square-applied}
L_{\mathrm{KL}}\bigl(\Phi(\gamma_nz)\bigr)=B\,L(\gamma_nz)=B\,L(z).
\end{equation}
By clause (1) of Definition~\ref{def:whole-tile-comparison}, after disk uniformization \(\Phi\) is a disk automorphism and extends to a homeomorphism \(\overline\Phi\) of the conformal ideal circles; the named-end compatibility \eqref{eq:named-end-compat} together with \eqref{eq:all-z} gives
\begin{equation}\label{eq:phi-limit}
\Phi(\gamma_nz)\longrightarrow\omega_+.
\end{equation}
The source-locked continuity of \(L_{\mathrm{KL}}\) at \(\omega_+\) (Definition~\ref{def:kl-tile-system}, \eqref{eq:kl-boundary-values}) now gives
\begin{equation}\label{eq:constant-value}
\overline L_{\mathrm{KL}}(\omega_+)=\lim_{n\to\infty}L_{\mathrm{KL}}\bigl(\Phi(\gamma_nz)\bigr)=B\,L(z).
\end{equation}
Equation \eqref{eq:constant-value} holds for \emph{every} \(z\in\mathcal E\).  The map \(L\) is nonconstant and \(B\) is injective on \(L(\mathcal E)\), so \(B\,L\) is nonconstant.  It cannot equal the one fixed value \(\overline L_{\mathrm{KL}}(\omega_+)\) for every \(z\).  This contradiction proves \eqref{eq:forced-trivial-kernel}.

The same proof works verbatim at \(\zeta_-\).  One wrapping vertex is sufficient; the second is an independent symmetric certificate.
\end{proof}